\documentclass[11pt]{article}
% Shared preamble for model_card.tex and assumptions_register.tex
\usepackage[T1]{fontenc}
\usepackage[utf8]{inputenc}
\usepackage{lmodern}
\usepackage[margin=2cm]{geometry}
\usepackage{amsmath,amssymb}
\usepackage{graphicx}
\usepackage{booktabs,longtable,tabularx,array,multirow}
\usepackage[table]{xcolor}
\usepackage{enumitem}
\usepackage{caption}
\usepackage[hidelinks]{hyperref}
\usepackage{xurl}
\usepackage{pdflscape}

\graphicspath{{figures/}}
\setlist{itemsep=2pt,topsep=3pt}
\captionsetup{font=small,labelfont=bf}
\renewcommand{\arraystretch}{1.15}

% ---- editing helpers -------------------------------------------------------
% \TODO{...} marks text the students still have to write or confirm (renders red in the PDF).
\newcommand{\TODO}[1]{\textcolor{red!75!black}{\textbf{[TODO:} #1\textbf{]}}}
% Single source for facts that appear in both documents; edit here, not in the text.
\newcommand{\Nall}{138{,}554}
\newcommand{\Nhh}{85{,}182}
\newcommand{\prev}{0.631}
\newcommand{\capacity}{30\%}
\newcommand{\epsTol}{0.10}
\newcommand{\minSupport}{30}
\newcommand{\nbDate}{25 October 2026}
% Team facts used by every document in writeup/, report/ and slides/. Edit here once.
\newcommand{\StudentA}{Sumedh Navuda}
\newcommand{\StudentB}{Aldino}
\newcommand{\Authors}{\StudentA{} and \StudentB{}}
\newcommand{\Contact}{sumedhnavuda@outlook.com}
\newcommand{\RollA}{[roll no.]}   % <- replace with Sumedh's roll number
\newcommand{\RollB}{[roll no.]}   % <- replace with Aldino's roll number
\newcommand{\RepoURL}{https://github.com/sumedhnvda/ict5143-p8-fairness-audits}
\newcommand{\TeamID}{[team ID]}   % <- replace with the team ID given by the instructor
\newcommand{\SubmitTag}{submission-v1}
   % names, roll numbers, contact, repository (shared with report/ and slides/)
\newcommand{\nb}{\texttt{Problem8\_\allowbreak Intersectional\_\allowbreak Fairness.ipynb}}
\newcommand{\sect}[1]{notebook~\S#1}
\newcommand{\status}[1]{\textsf{\small #1}}

\title{Individual Contribution Statements\\[2pt]\large ICT 5143 Capstone, Problem 8}
\date{\nbDate}
\author{\Authors}

\begin{document}\sloppy
\maketitle
\noindent The work was divided by pipeline stage. \StudentA{} led the question, literature, data, predictive Bayesian
network, exact inference and the Texas test of H1. \StudentB{} led the Bayesian audit models, the simulations, threshold
selection, uncertainty analysis and the improved audit procedure (H2). Each student reviewed the other's stages and can
explain the whole pipeline. Section numbers refer to \nb{}.

\section*{Contribution table}
\small
\begin{longtable}{@{}>{\raggedright\arraybackslash}p{4.0cm}>{\raggedright\arraybackslash}p{1.6cm}>{\raggedright\arraybackslash}p{1.6cm}>{\raggedright\arraybackslash}p{3.1cm}>{\raggedright\arraybackslash}p{4.3cm}@{}}
\toprule
Work item & Lead & Reviewer & Artefact & What the reviewer verified\\
\midrule\endhead
Research question, preregistration ($\epsilon$, capacity, utility budget, confirmatory family) & \StudentA & \StudentB & notebook top cells, \nolinkurl{docs/proposal.md} & confirmatory family matches \S11 Holm family\\
Literature review (19 papers), comparison table, narrowed claim, H1 & \StudentA & \StudentB & \nolinkurl{writeup/literature_review.tex} & closest-work rows against the struct and contamination audits\\
Data acquisition, label reversal, data card, household split & \StudentA & \StudentB & \S1--\S2 & no SERIALNO in two splits; label is exactly $1-$PUBCOV$_{=1}$\\
Predictive BN, VE, clique-tree check, LR/GBT baselines & \StudentA & \StudentB & \S3--\S4.1 & VE vs clique tree $3.3\times10^{-16}$; CI check\\
Experiment 1 (single vs intersectional) & \StudentA & \StudentB & \S5, \nolinkurl{exp1_*.csv} & 21 of 35 test groups with $\ge30$ positives\\
Audit models: EB, full Bayes, struct, contamination prior; verification & \StudentB & \StudentA & \S6, \S6.1 & beta-binomial posterior mean by hand; quadrature and gradient checks\\
Experiments 3 and 4, calibration, semi-synthetic worlds, robustness & \StudentB & \StudentA & \S7--\S9.2, \S12 & frozen \S9 rule reused unchanged in \S9.2\\
Experiments 2 and 5, bootstrap, Holm, survey weights & \StudentB & \StudentA & \S10, \S11, \S13 & Holm $p$ recomputed (0.368 $\to$ 0.736)\\
Texas replication and H1 test & \StudentA & \StudentB & \S14, \nolinkurl{h1_texas_test.csv} & decision rule fixed before the run; 27/27 count\\
Improved audit (design-matched calibration, held-out certification), H2 & \StudentB & \StudentA & \S15, \nolinkurl{v2_*.csv} & $J$ recomputed from \nolinkurl{v2_eval_*.csv}; paired CI\\
Manuscript, report, slides & both & both & \nolinkurl{writeup/}, \nolinkurl{report/}, \nolinkurl{slides/} & every number traced to \nolinkurl{results/}\\
Model card, assumptions register & \StudentB & \StudentA & \nolinkurl{writeup/} & statuses match notebook evidence\\
README, repository, reproducibility & \StudentA & \StudentB & \texttt{README.md}, \nolinkurl{docs/} & setup commands and manifest\\
\bottomrule
\end{longtable}
\normalsize

\section*{\StudentA{} (\RollA)}
\begin{description}[leftmargin=4.2cm,style=nextline]
  \item[Research design] Research question, preregistration table (30\% capacity, $\epsilon=0.10$, recall budget 0.02,
    confirmatory family) and the falsifiable contribution H1 with its Texas decision rule (\S14.1).
  \item[Literature review] 19 papers, the landscape comparison table and gap verification; narrowed the claim to an
    evaluation after finding that hierarchical intersectional estimation (Foulds et al.) and contamination priors
    (Berger and Berliner) already exist.
  \item[Data and predictive BN] Folktables acquisition and checks, data card, household-disjoint splits, discretisation,
    the factor and variable-elimination code, the BN, the inference verification against pgmpy and the LR/GBT baselines.
  \item[Experiments] Experiment 1; Texas replication and the H1 test (27 of 27 pooled claims failed on test).
  \item[Writing] Manuscript introduction, data and model sections; literature review; report sections 1--2 and
    reproducibility; slides 1--5; README and repository.
  \item[Technical decision] Splitting by household (SERIALNO) rather than by person, so that people who share
    circumstances never appear on both sides of a split.
  \item[Difficulty or failure] The earlier draft claim (the contamination prior halves false certification) failed:
    0.503 against the required 0.490. We kept the failure on record, narrowed the claim, and tested H1 out of sample.
  \item[Cross-check of \StudentB's work] Recomputed the beta-binomial posterior mean by hand ($k=2$, $n=10$, $\mu=0.4$,
    $\kappa=90$ gives 0.38), recomputed the Holm adjustment, and recomputed $J$ and its paired interval from
    \nolinkurl{results/v2_eval_*.csv}.
  \item[Tools and assistance] Python (numpy, scipy, pandas, pgmpy, folktables), Jupyter, LaTeX. An AI coding assistant
    (Claude Code) helped draft code, LaTeX and documentation; I checked every reported number against the saved notebook
    outputs in \nolinkurl{results/}.
  \item[Signature and date] \rule{6cm}{0.4pt}
\end{description}

\section*{\StudentB{} (\RollB)}
\begin{description}[leftmargin=4.2cm,style=nextline]
  \item[Audit models] Beta\nobreakdash-binomial empirical Bayes, the full-Bayes grid posterior, the attribute-structured audit
    (\texttt{struct}: marginal likelihood by adaptive quadrature with an exact gradient) and the contamination prior;
    their numerical and coverage checks (\S6.1).
  \item[Experiments] Experiment 3 on real counts; Experiment 4 (synthetic and semi-synthetic worlds) and the
    preregistered calibration (\S9); Experiments 2 and 5 (threshold family, minority thinning, selection); the household
    bootstrap with Holm correction; survey-weighted estimates; robustness grids.
  \item[Improved procedure] Diagnosis of the false alarms (90\% come from groups with fewer than 10 positives),
    design-matched calibration, held-out certification and the H2a--c tests (\S15).
  \item[Writing] Manuscript audit, experiment and improvement sections; model card; assumptions register; report
    sections 3--5; slides 6--10.
  \item[Technical decision] Ranking audit procedures by $J=\text{ACC}-\text{FTS}$ fixed before calibration, so that a
    procedure cannot win by never certifying (score 0) or always certifying (score 0).
  \item[Difficulty or failure] The calibration chosen on invented group sizes did not transfer to the real design: fair
    worlds were accepted 83\% of the time on the invented design and only 32--33\% on the California design. This
    motivated calibrating on the target design.
  \item[Cross-check of \StudentA's work] Checked that no household appears in two splits, that the label is the exact
    reversal of the Folktables label, that VE, pgmpy VE and the clique tree agree with enumeration, and the group
    denominators of Experiment 1.
  \item[Tools and assistance] Python (numpy, scipy, pandas), Jupyter, LaTeX. An AI coding assistant (Claude Code)
    helped draft code, LaTeX and documentation; I checked every reported number against the saved notebook outputs in
    \nolinkurl{results/}.
  \item[Signature and date] \rule{6cm}{0.4pt}
\end{description}
\end{document}
